%\usepackage{etex} %% suppression le 04/01/2024 j'ai l'impression qu'il ne sert plus à rien

\usepackage[utf8]{inputenc} 
%\usepackage[latin1]{inputenc} 

%\usepackage[T1]{fontenc} 	%apparence classique
\usepackage{fourier} 		%apparence Apmep

\usepackage[lmargin=1.5cm,rmargin=1.5cm,tmargin=1cm,bmargin=1.5cm]{geometry}
\usepackage{pstricks,pst-all,pst-plot,pst-text,pst-tree,pst-eps,pst-fill,pst-node,pst-math,pst-func,pst-solides3d,pst-3dplot,pstricks-add}
\usepackage{tikz,tkz-tab,tkz-linknodes}
\usetikzlibrary{lindenmayersystems}
\usetikzlibrary{calc,3d,shapes,backgrounds}
\usepackage{amsmath,amsfonts,amssymb,amsthm,makeidx}
\usepackage{animate}
\usepackage{array}
\usepackage{calc}
\usepackage{cancel}
%\usepackage[DVIps, leftbars]{changebar}\setlength{\changebarsep}{5mm}
\usepackage{coollist}
\usepackage{cprotect}
\usepackage{diagbox} % remplace slashbox
	% \diagbox[innerwidth=1.9cm, innerleftsep=0.4cm, innerrightsep=0.4cm]{$n$}{$k$}
\usepackage[b]{esvect}
\usepackage{eurosym}
\usepackage{fancybox}
\usepackage{fancyhdr}
\usepackage{fp}
%\usepackage{xfp}
\usepackage{framed}
\usepackage{graphicx} 
\usepackage{hhline}
\usepackage{ifthen}
\usepackage{lastpage}
\usepackage{lipsum}
\usepackage{longtable}
\usepackage{needspace}
\usepackage[version=3]{mhchem}
\usepackage{multicol}
\usepackage{multido}
\usepackage{multirow}
\usepackage[np]{numprint}
\usepackage{qrcode}
\usepackage{scratch3}
\usepackage{setspace}
\usepackage{siunitx}
%\sisetup{load-configurations = abbreviations}
\sisetup{locale = FR,detect-all,}
\usepackage{systeme}
\usepackage{tabularx}
\usepackage{tabulary}
\usepackage{tabto}
\usepackage{upgreek}
\usepackage{wrapfig}
\usepackage{xcolor,colortbl}

\pagestyle{fancy}
\fancyhf{}

\renewcommand\headrulewidth{0pt}

\setlength{\parindent}{0mm}
\definecolor{xlightgray}{gray}{0.9}

\newcommand*{\vect}[1]{\vv{\vphantom{k}#1}}
\newcommand*{\scal}[2]{\vect{#1}\cdot{}\vect{#2}}
\newcommand*{\veci}{\vect{\imath}}
\newcommand*{\vecj}{\vect{\jmath}}
\newcommand*{\vecn}{\vect{n}}
\newcommand*{\veck}{\vect{k}}
\newcommand*{\vecu}{\vect{u}}
\newcommand*{\vecv}{\vect{v}}
\newcommand*{\vecw}{\vect{w}}
\newcommand*{\veca}{\vect{a}}
\newcommand*{\vecb}{\vect{b}}
\newcommand*{\vecnul}{\vect{0}}
\newcommand*{\barre}[1]{\overline{\mathstrut#1\,}}
\newcommand*{\abs}[1]{\left\lvert\mathstrut#1\right\rvert}
\newcommand*{\norme}[1]{\lVert#1\rVert}
\newcommand*{\bignorme}[1]{\big\lVert#1\big\rVert}
\newcommand*{\lrnorme}[1]{\left\lVert\mathstrut#1\right\rVert}
\newcommand*{\intoo}[2]{\left] \, #1 \, ; \, #2 \, \right[}
\newcommand*{\intff}[2]{\left[ \, #1 \, ; \, #2 \, \right]}
\newcommand*{\intof}[2]{\left] \, #1 \, ; \, #2 \, \right]}
\newcommand*{\intfo}[2]{\left[ \, #1 \, ; \, #2 \, \right[}
\newcommand*{\intentier}[2]{\llbracket \, #1 \, ; \, #2 \, \rrbracket}
\newcommand*{\couple}[2]{\left( \, #1 \, ; \, #2 \, \right)}
\newcommand*{\triplet}[3]{\left( \, #1 \, ; \, #2 \, ; \, #3 \, \right)}
\newcommand*{\singleton}[1]{\left\lbrace \; #1 \; \right\rbrace}
\newcommand*{\paire}[2]{\left\lbrace \; #1 \; ; \; #2 \; \right\rbrace}
\newcommand*{\ensemble}[1]{%
	\begingroup%
		\newcommand{\pv}{\ensuremath{\; ; \;}}\left\lbrace \; #1 \;\right\rbrace%
	\endgroup}
	%% $A=\ensemble{1\pv 2\pv 3}$
\newcommand*{\uplet}[1]{%
	\begingroup%
		\newcommand{\pv}{\ensuremath{\, ; \,}}\left( \, #1 \,\right)%
	\endgroup}
\newcommand*{\anglevec}[2]{( \, \vect{#1} \, ; \, \vect{#2} \, )}
\newcommand*{\determinant}[2]{\det( \, \vect{#1} \, ; \, \vect{#2} \, )}
\newcommand*{\ii}{\text{i}}
\newcommand*{\jj}{\text{j}}
\newcommand*{\re}{\text{Re}}
\newcommand*{\im}{\text{Im}}
\newcommand*{\R}{\mathbb{R}}
\newcommand*{\N}{\mathbb{N}}
\newcommand*{\D}{\mathbb{D}}
\newcommand*{\Z}{\mathbb{Z}}
\newcommand*{\Q}{\mathbb{Q}}
\newcommand*{\C}{\mathbb{C}}
\newcommand*{\U}{\mathbb{U}}
\newcommand*{\cA}{\mathcal{A}}
\newcommand*{\cB}{\mathcal{B}}
\newcommand*{\cC}{\mathcal{C}}
\newcommand*{\cD}{\mathcal{D}}
\newcommand*{\cE}{\mathcal{E}}
\newcommand*{\cF}{\mathcal{F}}
\newcommand*{\cG}{\mathcal{G}}
\newcommand*{\cH}{\mathcal{H}}
\newcommand*{\cI}{\mathcal{I}}
\newcommand*{\cJ}{\mathcal{J}}
\newcommand*{\cK}{\mathcal{K}}
\newcommand*{\cL}{\mathcal{L}}
\newcommand*{\cM}{\mathcal{M}}
\newcommand*{\cN}{\mathcal{N}}
\newcommand*{\cO}{\mathcal{O}}
\newcommand*{\cP}{\mathcal{P}}
\newcommand*{\cQ}{\mathcal{Q}}
\newcommand*{\cR}{\mathcal{R}}
\newcommand*{\cS}{\mathcal{S}}
\newcommand*{\cT}{\mathcal{T}}
\newcommand*{\cU}{\mathcal{U}}
\newcommand*{\cV}{\mathcal{V}}
\newcommand*{\cW}{\mathcal{W}}
\newcommand*{\cX}{\mathcal{X}}
\newcommand*{\cY}{\mathcal{Y}}
\newcommand*{\cZ}{\mathcal{Z}}
\newcommand*{\Abarre}{\overline{A}}
\newcommand*{\Bbarre}{\overline{B}}
\newcommand*{\Cbarre}{\overline{C}}
\newcommand*{\Dbarre}{\overline{D}}
\newcommand*{\Ebarre}{\overline{E}}
\newcommand*{\Fbarre}{\overline{F}}
\newcommand*{\Gbarre}{\overline{G}}
\newcommand*{\Hbarre}{\overline{H}}
\newcommand*{\Ibarre}{\overline{I}}
\newcommand*{\Jbarre}{\overline{J}}
\newcommand*{\Kbarre}{\overline{K}}
\newcommand*{\Lbarre}{\overline{L}}
\newcommand*{\Mbarre}{\overline{M}}
\newcommand*{\Nbarre}{\overline{N}}
\newcommand*{\Obarre}{\overline{O}}
\newcommand*{\Pbarre}{\overline{P}}
\newcommand*{\Qbarre}{\overline{Q}}
\newcommand*{\Rbarre}{\overline{R}}
\newcommand*{\Sbarre}{\overline{S}}
\newcommand*{\Tbarre}{\overline{T}}
\newcommand*{\Ubarre}{\overline{U}}
\newcommand*{\Vbarre}{\overline{V}}
\newcommand*{\Wbarre}{\overline{W}}
\newcommand*{\Xbarre}{\overline{X}}
\newcommand*{\Ybarre}{\overline{Y}}
\newcommand*{\Zbarre}{\overline{Z}}

\newcommand*{\baseij}{(\, \veci \, , \, \vecj \,)}
\newcommand*{\baseijk}{(\, \veci \, , \, \vecj \, , \, \veck \,)}
\newcommand*{\baseuv}{(\, \vecu \, , \, \vecv \,)}
\newcommand*{\baseuvw}{(\, \vecu \, , \, \vecv \, , \, \vecw \,)}
\newcommand*{\Oij} {(\, O \, ; \, \veci , \vecj \,)}
\newcommand*{\Oijk}{(\, O \, ; \, \veci , \vecj , \veck \,)}
\newcommand*{\Ouv} {(\, O \, ; \, \vecu , \vecv \, )}
\newcommand*{\OIJ}{\triplet{O}{I}{J}}
\newcommand*{\repereplan}[3]{(\, #1 \, ; \, \vect{#2} \, , \, \vect{#3} \,)}
\newcommand*{\repereespace}[4]{(\, #1 \, ; \, \vect{#2} \, , \, \vect{#3} \, , \, \vect{#4} \, )}
\newcommand*{\baseplan}[2]  {(\, \vect{#1} \, , \, \vect{#2} \,)}
\newcommand*{\baseespace}[3]{(\,\vect{#1} \, , \,\vect{#2} \, ,\,\vect{#3}\,)}

\renewcommand*{\alpha}{\upalpha}
\newcommand*{\e}{\textrm{e}}
\newcommand*\circled[1]{\tikz[baseline=(char.base)]{ \node[shape=circle,draw,inner sep=1pt] (char) {#1};}}
\newcommand*{\grilleencadree}[2]{
	\begin{tikzpicture}[line cap=round,line join=round,>=triangle 45,x=1cm,y=1cm]
	\draw[color=lightgray,dash pattern=on 3pt off 3pt, xstep=0.75cm,ystep=0.75cm](0,0) grid (#1*0.75,#2*0.75);
	\draw[color=black](0,0) rectangle (#1*0.75,#2*0.75);
	\clip(0,0) rectangle (#1*0.75,#2*0.75);
	\end{tikzpicture}}
\newcommand*{\dotrule}[1]{%
   \parbox[t]{#1}{\dotfill}}   
\newcommand*{\asterism}{%
	\raisebox{-.5ex}{\setlength{\tabcolsep}{-.3pt}
	\begin{tabular}{cc}
	\multicolumn2c *\\[-2ex]*&*%
	\end{tabular}}}	
\newcommand*\clavier[1]{%
	\tikz[align=center, baseline={([yshift=-0.3ex]key.base)}]
 	\node[draw,
 	fill=white,
% 	drop shadow={shadow xshift=0.2ex,shadow yshift=-0.2ex,fill=black,opacity=0.2},
 	rectangle,
 	rounded corners=1.5pt,
 	inner sep=1pt,
 	line width=0.5pt,
 	font=\scriptsize\sffamily,
 	minimum width=1.6em]
 	(key){~#1~\strut};}
\newcommand*{\epve}{\; ; \;}
\newcommand*{\qpvq}{\quad ; \quad}
\newcommand*{\vq}{, \quad}
\newcommand*{\qetq}{\quad\textrm{et}\quad}
\newcommand*{\qouq}{\quad\textrm{ou}\quad}
\newcommand*{\qdouq}{\quad\textrm{d'où}\quad}
\newcommand*{\qsoitq}{\quad\textrm{soit}\quad}
\newcommand*{\qdoncq}{\quad\textrm{donc}\quad}
\newcommand*{\qeqq}{\quad\Longleftrightarrow\quad}

\newcommand*{\card}[1]{%
	\textrm{Card}(#1)}
\newcommand*{\tq}{%
	\;\mid\;}

%%% Exemple, remarque, exercices et problèmes numérotés
\newcommand*{\xmpl}{\textbf{Exemple}}
\newcommand*{\xmpls}{\textbf{Exemples}}
\newcommand*{\rmq}{\textbf{Remarque}}
\newcommand*{\rmqs}{\textbf{Remarques}}
\newcommand*{\demo}{\textbf{Démonstration}}
\newcommand*{\demos}{\textbf{Démonstrations}}
\newcommand*{\demoprg}{\textbf{Démonstration (au programme)}}

\newcounter{numexo}
\makeatletter
\newcommand{\exo}[1][\@nil]{%
	\def\tmp{#1}%
	\ifx\tmp\@nnil
		{\addtocounter{numexo}{1}\par\Needspace{7\baselineskip}\medskip\noindent\textbf{\underline{Exercice \arabic{numexo}}}\par\nopagebreak\smallskip}
	\else
		{\par\Needspace{7\baselineskip}\medskip\noindent\textbf{\underline{\smash{Exercice #1}}}\par\nopagebreak\smallskip}
	\fi}
\makeatother

\newcommand{\deuxchiffres}[1]{\ifnum#1<10 0\fi #1}

\newcounter{numpbl}
\makeatletter
\newcommand{\pbl}[1][\@nil]{%
	\def\tmp{#1}%
	\ifx\tmp\@nnil
		{\addtocounter{numpbl}{1}\medskip\noindent\textbf{\underline{Problème \arabic{numpbl}}}}
	\else
		{\medskip\noindent\textbf{\underline{\smash{Problème #1}}}}
	\fi}
\makeatother

%%% Théorème, définition
\newtheoremstyle{break}
	{0pt}		%Space before theorem
	{0pt}		%Space after theorem
	{\normalfont}			%Body font
	{}			%Indent amount (empty = no indent, \parindent = para indent)
	{\bfseries}	%Thm head font   
	{ -- }		%Punctuation after thm head
	{0pt}			%Space after thm head (\newline = linebreak)
	{}
\theoremstyle{break}
\newtheorem*{fthme}{Théorème}
\newtheorem*{fthmeadmis}{Théorème (admis)}
\newtheorem*{fdfnt}{Définition}
\newtheorem*{fdfnts}{Définitions}
\newtheorem*{fppte}{Propriété}
\newtheorem*{fcorol}{Corollaire}
\newtheorem*{fppteadmise}{Propriété (admise)}
\newtheorem*{fpptes}{Propriétés}
\newtheorem*{flem}{Lemme}
\newtheorem*{fmeth}{Méthode}
\newtheorem*{fboa}{Boîte à outils}
\newenvironment{thme}{\begin{framed}\begin{fthme}}{\end{fthme}\end{framed}}
\newenvironment{thmeadmis}{\begin{framed}\begin{fthmeadmis}}{\end{fthmeadmis}\end{framed}}
\newenvironment{dfnt}{\begin{framed}\begin{fdfnt}}{\end{fdfnt}\end{framed}}
\newenvironment{dfnts}{\begin{framed}\begin{fdfnts}}{\end{fdfnts}\end{framed}}
\newenvironment{ppte}{\begin{framed}\begin{fppte}}{\end{fppte}\end{framed}}
\newenvironment{corol}{\begin{framed}\begin{fcorol}}{\end{fcorol}\end{framed}}
\newenvironment{ppteadmise}{\begin{framed}\begin{fppteadmise}}{\end{fppteadmise}\end{framed}}
\newenvironment{pptes}{\begin{framed}\begin{fpptes}}{\end{fpptes}\end{framed}}
\newenvironment{lem}{\begin{framed}\begin{flem}}{\end{flem}\end{framed}}
\newenvironment{meth}{\begin{framed}\begin{fmeth}}{\end{fmeth}\end{framed}}
\newenvironment{boa}{\begin{framed}\begin{fboa}}{\end{fboa}\end{framed}}

\usepackage[english,french]{babel}
\DecimalMathComma
\frenchbsetup{StandardLists = true}

\usepackage{titlesec}
\renewcommand{\thesection}{\Roman{section}.}
\renewcommand{\thesubsection}{\arabic{subsection}.}
\renewcommand{\thesubsubsection}{\alph{subsubsection}.}
\titlespacing*{\section}		{0cm}	{2.0em plus 1.0em minus 0.2em}	{1.0em plus .2em}
\titlespacing*{\subsection}		{0.6cm}	{1.5em plus 0.5em minus 0.2em}	{0.5em plus .2em}
\titlespacing*{\subsubsection}	{1.2cm}	{1.0em plus 0.5em minus 0.2em}	{0.5em plus .2em}

%\usepackage{enumerate}
%\renewcommand{\theenumi}{\textbf{\arabic{enumi}}}
%\renewcommand{\theenumii}{\textbf{\alph{enumii}}}\renewcommand{\theenumiii}{\textit{\roman{enumiii}}}
%\renewcommand{\labelenumi}{\theenumi.}
%\renewcommand{\labelenumii}{\theenumii.}
%\renewcommand{\labelenumiii}{(\theenumiii)}
%\setlength\leftmargini{0.6cm} 	% 1er  niveau
%\setlength\leftmarginii{0.6cm} 	% 2ème niveau
%\setlength\leftmarginiii{0.6cm} % 3ème niveau
%\setlength\leftmarginiv{0.6cm} 	% 4ème niveau

\usepackage{titletoc}
\dottedcontents{section}%
  [15mm]		% espacement bord de page-titre
  {\bfseries}	
  {9mm}			% label width
  {3mm}			% leader width - espacement des points
\dottedcontents{subsection}%
  [21mm]
  {}
  {6mm}
  {3mm}

\usepackage{enumitem} % doit rester à la fin
\setlist[itemize]{
	topsep=-\parskip,
	partopsep=\parskip,
	parsep=0pt,
	itemsep=0pt,
	leftmargin=6mm,
	itemindent=0pt}	% OK
\setlist[enumerate]{
	topsep=-\parskip,
	partopsep=\parskip,
	parsep=0pt,
	itemsep=1.5pt, %	itemsep=0pt,
	labelwidth=0pt,
	leftmargin=6mm,
	itemindent=0pt} % OK
\setlist[enumerate,1]{
	label=\textbf{\arabic*.}}
\setlist[enumerate,2]{
	label=\textbf{\alph*.}}
\setlist[enumerate,3]{
	label=(\textit{\roman*\,})}

\newcommand{\titre}{
\boxput*(0.75,-1)
{\setlength{\fboxsep}{3pt}
\mbox{\includegraphics[height=0.8cm]{\logo}}} {%
\boxput*(-0.8,-1)
{\setlength{\fboxsep}{3pt}
\fcolorbox{white}{white}{\begin{LARGE}\type~\textbf{\num}\end{LARGE}}}
{\setlength{\fboxsep}{8pt}
\fcolorbox{black}{white}{
\begin{minipage}{0.97\linewidth}
\begin{center}\begin{LARGE}\textbf{\nom}\end{LARGE}\end{center}
\end{minipage}}
}}\vspace{1.6cm}}

\usepackage[colorlinks=true,linkcolor=blue,urlcolor=blue]{hyperref}

\newcommand*{\ds}{\displaystyle}

\newlength{\lgChaine}
\newcommand*{\tabcentretexte}[1]{%
	\settowidth{\lgChaine}{#1}%
	\tabto{\dimexpr 0.5\textwidth - 0.5\lgChaine} #1}
\newcommand*{\tabcentrepage}[1]{%
	\settowidth{\lgChaine}{#1}%
	\tabto{\dimexpr 0.5\linewidth - 0.5\lgChaine} #1}
